% ==================== 第 5 章 重点难点、创新点与可行性 ====================
\clearpage
\section{重点难点、创新点与可行性}

\subsection{研究重点与难点}

\subsubsection{业务规则的形式化与一致性}

同一条业务规则可能同时出现在数据录入、求解模型、校验逻辑和页面提示中。
如果各层口径不一致，系统即使“求解成功”也可能无法执行。本课题的重点是
建立唯一的数据字典和约束清单，并让求解器与独立校验器基于同一业务定义、
采用相互独立的实现路径。

\subsubsection{多目标权重与方案解释}

多目标加权能够生成差异化方案，但权重并不天然具有业务含义。研究难点在于
保证指标量纲统一、权重变化能够反映策略意图，并避免用一个总分掩盖局部
退化。课题将保留学生偏好、教师偏好、资源适配、负载均衡和变更数量等分项
指标，并开展权重敏感性实验。

\subsubsection{规模扩展与动态重排}

排课搜索空间随课次、教师、教室和时间槽数量快速增长。课题将通过不可行
组合预筛选、对称性削减、求解时限和分区求解等方法控制规模。调课场景还要
在新增约束与方案稳定之间权衡，因此需要对“最少变更”建立清晰的度量和
验收方法。

\subsection{拟创新点}

\begin{enumerate}
  \item \textbf{求解与独立校验的双层保障。}求解器负责生成候选课表，
  校验器直接基于业务对象复核结果，使“可行”结论不依赖单一代码路径。
  \item \textbf{面向教务决策的多策略解释。}把不同角色关注点拆为统一分项
  指标，通过策略权重生成候选方案，并解释每个方案的得失，而不是只展示
  一个不可解释的总目标值。
  \item \textbf{基于基线差异的最小影响调课。}把原课表作为显式输入，以
  变更课次为惩罚项，输出新方案的同时生成逐条差异和受影响对象。
  \item \textbf{算法核心与协同平台解耦。}以稳定 API 封装排课能力，再接入
  多维表格、消息和审批，使求解逻辑可以独立测试，也便于替换交互端。
\end{enumerate}

上述创新点是面向具体应用的系统与方法创新，不声称对通用 CP-SAT 求解理论
作出突破。其有效性需要通过后续实验和用户评价验证。

\subsection{技术可行性}

OR-Tools CP-SAT 已提供布尔变量、整数变量、资源互斥和线性目标等建模能力，
能够直接表达本课题的核心约束。Python 适合快速实现数据模型、校验器和服务
接口；原生 Web 技术可满足原型展示；自动化测试则为持续修改约束提供回归
保障。当前可运行原型进一步证明了基础技术链路可打通。

\subsection{经济与实施可行性}

课题开发阶段以开源软件和本地运行环境为主，不依赖专用硬件。模块化设计使
算法、接口和前端可以逐步迭代；飞书集成放在核心求解稳定之后，可降低平台
权限、回调配置和接口变化对算法研发的干扰。实际部署成本将在真实规模实验
后，结合算力、运维和接口调用量重新评估。

\subsection{合规可行性}

演示数据使用虚构身份和资源信息。进入真实试点前，需要落实最小权限、访问
控制、传输加密、日志脱敏、数据保存期限和删除机制；飞书应用凭据不得进入
前端或源码仓库。若涉及未成年人信息，还应在组织授权和监护人知情范围内
处理，并仅采集完成排课所必需的字段。
